% Hand-written discussion fragment for the drive-compliance map.
% NOT generated by maestro2tex -- the tool only places it, via the "order" list in
% configs/anatop_tb_controlamp.json.
%
% Every number here comes from tab_ctrlamp_drivemap.tex (typical corner),
% fig_ctrlamp_drivespan.tex (all five), tab_ctrlamp_ac.tex and
% tab_ctrlamp_accuracy.tex, and is regenerated with them. Re-read this against
% those tables after any re-run.

\paragraph{Reading the map.}
The compliance envelope is not widest at mid-supply. It is pinched to under
\SI{4}{\micro\ampere} of span below \SI{0.5}{\volt}, widens steadily through the middle
of the range to a maximum of \SI{13.64}{\micro\ampere} at \SI{1.8}{\volt}
(\SIrange{-6.04}{7.60}{\micro\ampere}), then collapses above \SI{1.9}{\volt}. Closed-loop
output resistance follows the same shape inverted, from \SI{20.7}{\kilo\ohm} at the
bottom of the range down to \SI{5.6}{\kilo\ohm} at \SI{1.7}{\volt} and back to
\SI{283}{\kilo\ohm} at \SI{2.0}{\volt}: where the output devices sit deep in saturation
the loop holds the node stiffly, and where they do not, it cannot. At the mid-supply
point the potentiostat would actually sit on, \SI{1.2}{\volt}, the usable range is
\SIrange{-3.49}{3.36}{\micro\ampere} on \SI{55.18}{\micro\ampere} of supply.

Two features are worth reading off the map rather than the table. First, compliance is
close to \emph{symmetric} in the middle of the range --- the sinking and sourcing edges
stay within about \SI{20}{\percent} of each other from \SI{0.3}{\volt} to
\SI{1.7}{\volt} --- which is nothing like the four-to-one sourcing advantage the slew
bench shows. Second, below \SI{0.3}{\volt} and above \SI{2.1}{\volt} the zero-load error
column already exceeds \SI{50}{\milli\volt}: the buffer does not track its input at any
load there, so those rows describe the external current needed to drag the output onto
the commanded value, not a drive capability. The envelope legitimately ceases to exist
outside the tracking range, which is a property of the block worth seeing rather than a
gap in the data.

Against the rev-2 target of \(\pm\SI{33}{\micro\ampere}\) into the counter electrode, the
amplifier as biased is roughly ten times short at its best point and twenty times short
at the operating point. Figure~\ref{fig:ctrlamp-drivespan} shows the same envelope in
every corner: the shape is identical and the peak moves only between
\SI{12.40}{\micro\ampere} (ff) and \SI{14.76}{\micro\ampere} (ss). The deficit is
structural, not a corner artefact, and no trim code reaches the target --- raising the
bias lifts both edges together but does not change the ratio or close a factor of ten.

\paragraph{The ff corner behaves very differently, and that should be checked.}
In ff the amplifier keeps its bias --- \SI{31.51}{\micro\ampere} of supply against
\SI{31.81}{\micro\ampere} typical, with all four cascode rails in family --- but its
small-signal behaviour collapses: open-loop gain falls to \SI{72.67}{\decibel} from
\SI{83.87}{\decibel}, the dominant pole drops to \SI{20.1}{\hertz} from
\SI{128.9}{\hertz}, and gain--bandwidth falls from \SI{2.018}{\mega\hertz} to
\SI{87}{\kilo\hertz}, a factor of twenty-three. The usable unity-gain input window
shrinks from \SIrange{0.302}{2.006}{\volt} to \SIrange{0.953}{2.057}{\volt}, and the
loop-gain phase never reaches \SI{-180}{\degree}, which is why ff has no gain margin
entry in Table~\ref{tab:ctrlamp-stability}. Five independent measurements agree, and the
open-loop response is a clean single-pole roll-off, so this is not an extraction
artefact.

A likely explanation is the corner set itself rather than the silicon: fast transistors
with typical resistors is not a physically self-consistent corner, and this block's bias
current is set by a poly resistor. That combination can put the input stage near the
edge of its range while the total current still looks correct, which is what the
shrunken input window suggests. It is a hypothesis, not a result --- it should be
settled by re-running with the resistor rows tracking the transistor rows before any ff
number here is used for a design decision.

\paragraph{The corner rows are still pinned.}
The run covers five process corners, but the imported corner set selects
\texttt{tt\_res} from \texttt{cor\_res.scs} and \texttt{tt\_mim} from
\texttt{cor\_mim.scs} in every one of them; only \texttt{cor\_std\_mos.scs} actually
moves. A separate spot check at \SI{1.25}{\volt} gives \(\pm\SI{1.7}{\micro\ampere}\) for
ff transistors with \texttt{tt\_res} against \(\pm\SI{6.2}{\micro\ampere}\) with
\texttt{ff\_res}, on \SI{54.9}{\micro\ampere} of supply against
\SI{85.5}{\micro\ampere} --- the resistor corner alone moves the drive by a factor of
\num{3.6}, more than any transistor corner in
Figure~\ref{fig:ctrlamp-drivespan} does. So the spread shown throughout this section is
the transistor-corner spread only. Fix the corner rows and re-run before treating any of
it as signed off.
